\chapter{Correlation, Transformation, and Scaling}
\label{chap:correlation-normalization}

\chapterguide
  {You should be comfortable loading data, selecting columns, using packages, and creating plots.}
  {compute and test Pearson correlation, visualize correlation matrices, and compare log transformation, standardization, and min--max scaling.}

\section{Pearson correlation with \texttt{cor()}}

Compute a Pearson coefficient for two paired vectors:

\begin{lstlisting}[style=dsr]
x <- c(1, 2, 3, 4, 5, 6, 7)
y <- c(1, 3, 6, 2, 7, 4, 5)

result <- cor(x, y, method = "pearson")
cat("Pearson correlation coefficient is:", result)
\end{lstlisting}

The coefficient summarizes their linear association from $-1$ to $1$. That
range is easier to read once you have seen what each value looks like.

\begin{utpfig}[What r looks like]{Forty points at five correlation strengths.
$r$ measures how tightly the cloud hugs a straight line, and its sign gives the
direction. It says nothing about the slope: a steep line and a shallow one can
both have $r = 0.95$.}{fig:corr-gallery}
% Generated by tools/make-figures.py -- do not edit by hand.
% Regenerate with: python3 tools/make-figures.py
\begin{tikzpicture}
\begin{groupplot}[
  group style={group size=5 by 1, horizontal sep=6mm},
  width=3.5cm, height=3.5cm,
  xmin=-2.6, xmax=2.6, ymin=-2.6, ymax=2.6,
  xtick=\empty, ytick=\empty,
  axis line style={draw=UTPMuted!70},
  title style={font=\scriptsize\bfseries\color{UTPNavy}},
]
\nextgroupplot[title={$r = -0.94$}]
\addplot+[only marks, mark=*, mark size=1.0pt, draw=UTPRed, fill=UTPRed]
  coordinates {(0.489,-0.416) (0.856,-0.765) (0.149,-0.598) (-0.611,0.564) (-0.073,0.548) (-0.736,0.082) (0.562,-0.208) (2.141,-1.903) (-0.120,-0.142) (-0.278,0.967) (1.092,-0.722) (0.928,-1.291) (0.618,-0.545) (-0.660,0.791) (0.452,-0.492) (1.346,-0.979) (-1.171,1.017) (-0.077,0.298) (-1.858,2.187) (-1.104,0.733) (-1.785,1.669) (0.198,-0.363) (-1.076,1.002) (0.822,-1.192) (0.681,-0.842) (0.257,-0.319) (-2.618,2.675) (-0.177,0.563) (0.428,-0.886) (0.627,-0.872) (-1.400,1.484) (-0.102,-0.653) (-0.720,0.464) (-0.510,0.409) (1.797,-1.176) (-0.509,0.696) (0.447,-0.539) (1.579,-1.575) (-0.232,0.103) (0.350,0.226)};
\nextgroupplot[title={$r = -0.42$}]
\addplot+[only marks, mark=*, mark size=1.0pt, draw=UTPRed, fill=UTPRed]
  coordinates {(-0.421,0.512) (-0.268,-0.224) (0.541,-0.132) (-0.042,-1.703) (-0.137,-0.812) (-1.267,1.167) (0.092,-1.892) (-0.440,1.261) (1.544,-2.183) (0.912,0.554) (0.077,-0.630) (0.931,0.567) (-0.312,0.503) (1.404,-1.914) (0.100,1.408) (0.826,1.263) (-1.485,0.847) (0.534,-0.281) (-1.975,0.980) (-2.403,0.378) (-0.269,1.429) (-1.003,0.160) (0.309,0.040) (2.874,-1.860) (-0.919,0.040) (-0.663,-0.714) (0.360,1.296) (0.714,-0.246) (-0.111,1.227) (-0.147,0.312) (0.973,-0.893) (0.748,-0.430) (-1.168,0.219) (0.009,0.235) (0.150,-0.204) (-1.194,0.485) (0.427,-1.312) (-0.951,-0.122) (1.305,1.252) (0.344,-0.582)};
\nextgroupplot[title={$r = -0.09$}]
\addplot+[only marks, mark=*, mark size=1.0pt, draw=UTPMuted, fill=UTPMuted]
  coordinates {(-1.296,-1.003) (0.388,1.082) (1.683,0.718) (-1.780,1.093) (-0.272,2.088) (-0.785,1.698) (-2.151,-0.903) (0.870,-1.441) (-0.048,1.037) (0.067,-0.774) (-0.930,0.217) (0.531,1.060) (-0.673,-1.396) (-0.193,-1.857) (-1.361,0.486) (-1.492,0.273) (1.597,-0.014) (-0.634,0.268) (0.333,-0.622) (-0.061,-1.267) (-0.554,0.065) (-0.635,-0.732) (0.743,-1.396) (-0.385,0.730) (-0.203,0.169) (0.007,0.631) (1.404,-0.749) (0.804,-0.421) (0.443,-0.759) (-0.702,-0.647) (-1.692,0.423) (1.129,-0.545) (1.150,0.582) (-0.190,0.565) (0.636,2.232) (0.926,-1.148) (0.986,1.108) (1.095,0.141) (-0.571,0.216) (1.814,-1.204)};
\nextgroupplot[title={$r = +0.65$}]
\addplot+[only marks, mark=*, mark size=1.0pt, draw=UTPBlue, fill=UTPBlue]
  coordinates {(-0.269,-0.142) (0.844,-0.579) (0.080,-0.269) (0.232,-0.095) (-0.112,0.987) (1.224,0.969) (0.137,0.012) (-1.878,0.272) (-0.483,-0.868) (-1.664,-1.404) (-2.849,-3.440) (-0.334,1.215) (1.389,1.657) (0.200,0.915) (-0.927,0.021) (-0.713,-0.378) (1.202,0.844) (0.302,-0.122) (0.201,-0.136) (0.107,0.051) (0.192,1.469) (0.901,0.987) (0.667,0.274) (0.356,-0.398) (-0.808,-1.133) (0.629,1.811) (-0.476,0.134) (1.168,0.685) (-1.019,-0.985) (0.029,-1.087) (0.150,-0.038) (-1.068,0.131) (1.749,0.594) (1.507,0.607) (-0.272,-0.430) (0.870,0.550) (0.507,-1.262) (-2.260,-1.609) (0.388,-0.638) (0.100,0.826)};
\nextgroupplot[title={$r = +0.94$}]
\addplot+[only marks, mark=*, mark size=1.0pt, draw=UTPBlue, fill=UTPBlue]
  coordinates {(-0.958,-0.667) (0.582,1.237) (1.664,1.795) (-0.436,-0.339) (-0.765,-1.251) (0.141,0.312) (-0.080,-0.706) (-1.212,-1.570) (0.449,0.760) (2.225,2.355) (-0.372,-0.352) (-0.309,-0.837) (-1.271,-1.263) (0.287,0.076) (-1.502,-1.126) (-1.342,-0.413) (1.384,1.405) (-1.112,-1.257) (1.158,0.712) (1.664,1.569) (0.219,0.187) (0.555,0.159) (2.153,2.083) (-0.302,-0.635) (-0.755,-0.267) (-1.623,-1.090) (-0.030,0.396) (1.612,1.376) (1.019,1.170) (-1.154,-1.070) (-1.055,-1.068) (-0.653,-0.678) (0.256,-0.169) (-0.312,-0.746) (0.168,0.478) (0.262,0.222) (-0.419,-0.267) (-0.123,0.280) (0.159,-0.411) (-0.173,-0.388)};
\end{groupplot}
\end{tikzpicture}

\end{utpfig}

\section{Pearson correlation with \texttt{cor.test()}}

\texttt{cor.test()} adds inferential output:

\begin{lstlisting}[style=dsr]
result <- cor.test(x, y, method = "pearson")
print(result)
\end{lstlisting}

Its output includes the estimate, test statistic, p-value, and confidence interval.

\begin{pitfall}
Correlation measures association, not causation. Inspect the data and assumptions before interpreting the coefficient or p-value.
\end{pitfall}

\subsection{Why the coefficient is not enough on its own}
Anscombe's quartet is four small datasets built so that the mean, variance,
regression line, and correlation are the same in all four. Only a plot
separates them.

\begin{utpfig}[Anscombe's quartet]{Four datasets with identical summary
statistics: $r = 0.82$ and the same fitted line in every panel. Set 1 is a
genuine linear relationship. Set 2 is a curve. Set 3 is a perfect line plus
one outlier. Set 4 has no relationship at all --- a single point at $x = 19$
manufactures the whole correlation.}{fig:anscombe}
% Generated by tools/make-figures.py -- do not edit by hand.
% Regenerate with: python3 tools/make-figures.py
\begin{tikzpicture}
\begin{groupplot}[
  group style={group size=4 by 1, horizontal sep=7mm},
  width=4.0cm, height=3.8cm,
  xmin=2, xmax=20, ymin=2, ymax=14,
  xtick={5,10,15}, ytick={4,8,12},
  tick label style={font=\tiny},
  axis line style={draw=UTPMuted!70},
  grid=major, grid style={draw=UTPMuted!15},
  title style={font=\scriptsize\bfseries\color{UTPNavy}},
]
\nextgroupplot[title={set 1: $r=0.82$}]
\addplot+[only marks, mark=*, mark size=1.3pt, draw=UTPBlue, fill=UTPBlue]
  coordinates {(10.00,8.04) (8.00,6.95) (13.00,7.58) (9.00,8.81) (11.00,8.33) (14.00,9.96) (6.00,7.24) (4.00,4.26) (12.00,10.84) (7.00,4.82) (5.00,5.68)};
\addplot+[domain=3:20, samples=2, draw=UTPRed, thick, solid, mark=none]
  {0.5001*x + 3.0001};
\nextgroupplot[title={set 2: $r=0.82$}]
\addplot+[only marks, mark=*, mark size=1.3pt, draw=UTPBlue, fill=UTPBlue]
  coordinates {(10.00,9.14) (8.00,8.14) (13.00,8.74) (9.00,8.77) (11.00,9.26) (14.00,8.10) (6.00,6.13) (4.00,3.10) (12.00,9.13) (7.00,7.26) (5.00,4.74)};
\addplot+[domain=3:20, samples=2, draw=UTPRed, thick, solid, mark=none]
  {0.5000*x + 3.0009};
\nextgroupplot[title={set 3: $r=0.82$}]
\addplot+[only marks, mark=*, mark size=1.3pt, draw=UTPBlue, fill=UTPBlue]
  coordinates {(10.00,7.46) (8.00,6.77) (13.00,12.74) (9.00,7.11) (11.00,7.81) (14.00,8.84) (6.00,6.08) (4.00,5.39) (12.00,8.15) (7.00,6.42) (5.00,5.73)};
\addplot+[domain=3:20, samples=2, draw=UTPRed, thick, solid, mark=none]
  {0.4997*x + 3.0025};
\nextgroupplot[title={set 4: $r=0.82$}]
\addplot+[only marks, mark=*, mark size=1.3pt, draw=UTPBlue, fill=UTPBlue]
  coordinates {(8.00,6.58) (8.00,5.76) (8.00,7.71) (8.00,8.84) (8.00,8.47) (8.00,7.04) (8.00,5.25) (19.00,12.50) (8.00,5.56) (8.00,7.91) (8.00,6.89)};
\addplot+[domain=3:20, samples=2, draw=UTPRed, thick, solid, mark=none]
  {0.4999*x + 3.0017};
\end{groupplot}
\end{tikzpicture}

\end{utpfig}

\begin{keyidea}
Plot before you report. \texttt{cor()} returns a number for any two numeric
vectors, including the three above where the number is misleading. A
scatterplot costs one line and rules out all three failures at once.
\end{keyidea}

\section{Correlation matrix with \texttt{corrplot}}

Compute and visualize the \texttt{mtcars} correlation matrix:

\begin{lstlisting}[style=dsr]
install.packages("corrplot")
install.packages("RColorBrewer")

library(corrplot)
library(RColorBrewer)

M <- cor(mtcars)
?corrplot

corrplot(M, type = "upper")
corrplot(M, type = "upper", order = "hclust")
corrplot(
  M,
  type = "upper",
  order = "hclust",
  col = brewer.pal(n = 8, name = "RdYlBu")
)
\end{lstlisting}

Hierarchical ordering places variables with similar correlation patterns together.

\section{Building a heatmap-ready correlation table}

Use the \texttt{environmental} dataset to prepare a clustered heatmap:

\begin{lstlisting}[style=dsr]
install.packages("lattice")
library(lattice)

data <- environmental
head(data)

corr_mat <- round(cor(data), 2)
head(corr_mat)
\end{lstlisting}

The matrix is reordered using a distance derived from correlation values and hierarchical clustering:

\begin{lstlisting}[style=dsr]
corr_dist <- as.dist((1 - corr_mat) / 2)
hc <- hclust(corr_dist)
corr_mat <- corr_mat[hc$order, hc$order]
\end{lstlisting}

Convert the matrix to long form for \texttt{ggplot2}:

\begin{lstlisting}[style=dsr]
install.packages("reshape2")
library(reshape2)

melted_corr_mat <- melt(corr_mat)
head(melted_corr_mat)
\end{lstlisting}


\section{Correlation heatmaps}

\subsection{Using \texttt{ggplot2}}

\begin{lstlisting}[style=dsr]
install.packages("ggplot2")
library(ggplot2)

ggplot(
  data = melted_corr_mat,
  aes(x = Var1, y = Var2, fill = value)
) +
  geom_tile() +
  geom_text(
    aes(label = value),
    color = "white",
    size = 4
  )
\end{lstlisting}

Tile color and overlaid text both encode the correlation value.

\subsection{Using \texttt{heatmaply}}

\begin{lstlisting}[style=dsr]
install.packages("heatmaply")
library(heatmaply)

heatmaply_cor(
  x = cor(data),
  xlab = "Features",
  ylab = "Features",
  k_col = 2,
  k_row = 2
)
\end{lstlisting}

\subsection{Using \texttt{ggcorrplot}}

Use the correctly named package:

\begin{lstlisting}[style=dsr]
install.packages("ggcorrplot")
library(ggcorrplot)
ggcorrplot::ggcorrplot(cor(data))
\end{lstlisting}

\begin{pitfall}
The source misspells the installation name as \texttt{ggcorplot}. The package and namespace are \texttt{ggcorrplot}.
\end{pitfall}

\section{Log transformation}

The supplied numeric vector is:

\begin{lstlisting}[style=dsr]
mydata <- c(244, 753, 596, 645, 874, 141, 639, 465, 999, 654)

scaled_data <- log(mydata)
print(scaled_data)
\end{lstlisting}

Log transformation compresses large values and can reduce right skew; it does not constrain values to a fixed range.

The three methods in this chapter are easiest to compare on the same ten
numbers, plotted side by side.

\begin{utpfig}[Three scalings compared]{The same ten values under each
transformation. The shape of the cloud is unchanged by \texttt{scale()} and
min--max --- only the axis labels move. \texttt{log()} is the one that alters
the spacing between points, which is why it is the only one of the three that
can change the outcome of a later correlation.}{fig:scaling}
% Generated by tools/make-figures.py -- do not edit by hand.
% Regenerate with: python3 tools/make-figures.py
\begin{tikzpicture}
\begin{groupplot}[
  group style={group size=4 by 1, horizontal sep=13mm},
  width=3.6cm, height=3.5cm,
  xmin=0.4, xmax=10.6, xtick=\empty,
  xlabel style={font=\tiny\color{UTPMuted}, at={(0.5,-0.02)}},
  scaled y ticks=false,
  tick label style={font=\tiny},
  axis line style={draw=UTPMuted!70},
  grid=major, grid style={draw=UTPMuted!15},
  title style={font=\scriptsize\bfseries\color{UTPNavy}},
]
\nextgroupplot[title={\texttt{raw}}, xlabel={\tiny 141 to 999}]
\addplot+[only marks, mark=*, mark size=1.4pt, draw=UTPMuted, fill=UTPMuted]
  coordinates {(1,244.0000) (2,753.0000) (3,596.0000) (4,645.0000) (5,874.0000) (6,141.0000) (7,639.0000) (8,465.0000) (9,999.0000) (10,654.0000)};
\nextgroupplot[title={\texttt{log()}}, xlabel={\tiny 4.9 to 6.9}]
\addplot+[only marks, mark=*, mark size=1.4pt, draw=UTPTeal, fill=UTPTeal]
  coordinates {(1,5.4972) (2,6.6241) (3,6.3902) (4,6.4693) (5,6.7731) (6,4.9488) (7,6.4599) (8,6.1420) (9,6.9068) (10,6.4831)};
\nextgroupplot[title={\texttt{scale()}}, xlabel={\tiny mean 0, sd 1}]
\addplot+[only marks, mark=*, mark size=1.4pt, draw=UTPBlue, fill=UTPBlue]
  coordinates {(1,-1.3604) (2,0.5792) (3,-0.0191) (4,0.1677) (5,1.0403) (6,-1.7529) (7,0.1448) (8,-0.5182) (9,1.5166) (10,0.2020)};
\nextgroupplot[title={\texttt{min--max}}, xlabel={\tiny 0 to 1}]
\addplot+[only marks, mark=*, mark size=1.4pt, draw=UTPGreen, fill=UTPGreen]
  coordinates {(1,0.1200) (2,0.7133) (3,0.5303) (4,0.5874) (5,0.8543) (6,0.0000) (7,0.5804) (8,0.3776) (9,1.0000) (10,0.5979)};
\end{groupplot}
\end{tikzpicture}

\end{utpfig}

\begin{center}
\begin{tabularx}{\linewidth}{@{}p{2.5cm}p{2.9cm}X@{}}
\toprule
\rowcolor{DSPaleBlue!92}
Method & Output range & Use it when \\
\midrule
\texttt{log()} & unbounded & values span orders of magnitude or are right-skewed \\
\texttt{scale()} & mean 0, sd 1 & comparing variables measured in different units \\
min--max & 0 to 1 & an algorithm needs a bounded input \\
\bottomrule
\end{tabularx}
\end{center}

\section{Standardization}

\texttt{scale()} centers values at approximately zero and scales them to unit standard deviation:

\begin{lstlisting}[style=dsr]
scaled_data2 <- as.data.frame(scale(mydata))
print(scaled_data2)
\end{lstlisting}


\section{Min--max scaling}

Min--max scaling maps values to a fixed range, normally 0 to 1:

\begin{lstlisting}[style=dsr]
install.packages("caret")
library(caret)

minmax <- preProcess(
  as.data.frame(mydata),
  method = c("range")
)

scaled_data3 <- predict(
  minmax,
  as.data.frame(mydata)
)
print(scaled_data3)
\end{lstlisting}

\section{Comparing summaries}

Compare the raw and transformed summaries:

\begin{lstlisting}[style=dsr]
summary(mydata)
summary(scaled_data)
summary(scaled_data2)
summary(scaled_data3)
\end{lstlisting}

\begin{pitfall}
The source calls an undefined object, \texttt{scaled\_data1}. The corrected comparison uses \texttt{scaled\_data}.
\end{pitfall}

\begin{dsfigure}
\begin{tikzpicture}[
  box/.style={draw=DSBlue,fill=DSPaleBlue,rounded corners=2mm,text width=3.35cm,minimum height=1.1cm,align=center,font=\small\bfseries\color{DSNavy}},
  arr/.style={-{Stealth[length=2mm]},thick,draw=DSTeal}
]
\node[box] (raw) {Raw numeric data};
\node[box,right=1.0cm of raw] (log) {Log transformation};
\node[box,below=0.8cm of log] (std) {Standard scaling};
\node[box,above=0.8cm of log] (mm) {Min--max scaling};
\draw[arr] (raw) -- (log);
\draw[arr] (raw.east) -- ++(0.55,0) |- (std.west);
\draw[arr] (raw.east) -- ++(0.55,0) |- (mm.west);
\end{tikzpicture}
\end{dsfigure}

Compare each method with the raw data because the resulting scales have different interpretations.

\begin{quickcheck}
\begin{enumerate}
  \item What is the difference in output intent between \texttt{cor()} and \texttt{cor.test()} in Lab 9a?
  \item Which dataset is used for the first \texttt{corrplot()} example?
  \item Why is \texttt{melt()} applied to the reordered correlation matrix before the \texttt{ggplot2} heatmap?
  \item Which function performs the standard scaling in the handout?
  \item Which object-name mismatch appears in the final summary comparison?
\end{enumerate}
\end{quickcheck}

\begin{chapterrecap}
\begin{itemize}
  \item \texttt{cor()} estimates association; \texttt{cor.test()} adds inferential statistics.
  \item Clustered heatmaps group variables with similar correlation patterns.
  \item Log transformation, standardization, and min--max scaling serve different purposes.
\end{itemize}
\end{chapterrecap}
